\chapter{Implementation and Development}
\chaptertoc

\section{Overview}
The project is structured in five Scrum sprints, each delivering an operational increment of the final system, from core computer vision to uncertainty-aware monitoring and external integration.

\section{Sprint I: Core Detection and Queue Logic}

\subsection{Sprint Backlog}
Main backlog items for this sprint were person detection, multi-object tracking, and queue-zone filtering.

\subsection{System Architecture}
The foundation layer consists of three core components:
\begin{itemize}
    \item YOLO26 person detection\cite{ultralytics_docs},
    \item ByteTrack object tracking\cite{bytetrack2022},
    \item Polygon zone filtering.
\end{itemize}

\begin{figure}[ht]
    \centering
    \fbox{\parbox[c][6cm][c]{0.8\textwidth}{\centering Placeholder for architecture diagram (add architecture_diagram.png)}}
    \caption{High-level system architecture: detectors, tracker, zone manager, analyzer, and integration components.}
    \label{fig:architecture}
\end{figure}

See Figure~\ref{fig:architecture} for a high-level overview of the main components and data flow.

\subsection{Module Implementation}
\subsubsection{Detector Module}
The detector module wraps Ultralytics YOLO for real-time person detection:
\begin{itemize}
    \item Supports model sizes: nano, small, medium, large, xlarge,
    \item Returns bounding boxes with confidence scores.
\end{itemize}

\subsubsection{Tracker Module}
ByteTrack assigns consistent IDs across frames to prevent duplicate counting.

\subsubsection{Zone Manager}
Polygon-based filtering ensures only people inside the configured queue zone are analyzed.

\section{Sprint II: Graphical Interface and Multi-Video Support}

\subsection{Sprint Backlog}
Main backlog items for this sprint were user workflow design, per-video zone configuration, and multi-source execution.

\subsection{GUI Development}
The interface follows a simple workflow:
\begin{enumerate}
    \item choose the number of videos,
    \item select sources,
    \item configure the model,
    \item define one zone per video,
    \item launch parallel analysis.
\end{enumerate}

\subsection{Interface Technologies}
\subsubsection{Tkinter and CustomTkinter}
Primary framework for the graphical interface with custom styling.

\subsubsection{Zone Selector Tool}
Interactive polygon selection for queue area definition per video source.

\section{Sprint III: YOLO26 Migration and Performance Tuning}

\subsection{Sprint Backlog}
Main backlog items for this sprint were model migration, inference optimization, and runtime stability improvements.

\subsection{Model Selection}
Migration from earlier YOLO versions to YOLO26 nano/small for optimal speed-accuracy trade-off.

\subsection{Performance Optimization}
\subsubsection{GPU Acceleration}
PyTorch CUDA support for real-time inference.

\subsubsection{Inference Settings}
Tuning confidence thresholds and NMS parameters for operational requirements.

\section{Sprint IV: Uncertainty Quantification}

\subsection{Sprint Backlog}
Main backlog items for this sprint were Bayesian rate modeling, confidence intervals, and uncertainty-level classification.

\subsection{Bayesian Rate Estimation}
Rates are modeled using posterior Gamma distributions:
\begin{equation}
\lambda \sim \Gamma(\alpha_{\lambda}, \beta_{\lambda}),
\qquad
\mu \sim \Gamma(\alpha_{\mu}, \beta_{\mu}).
\end{equation}

This Bayesian approach follows standard inference patterns in machine learning and statistics \cite{bishop_prml},
and is combined with queueing-theory reasoning for interpretability \cite{gross_queueing}.

\subsection{Confidence Intervals}
Ninety-five percent credible intervals are computed for each metric.

\subsection{Uncertainty Levels}
Classification of uncertainty as Low, Medium, or High based on posterior variance.

\section{Sprint V: n8n Integration and Alert System}

\subsection{Sprint Backlog}
Main backlog items for this sprint were webhook integration, payload standardization, alert threshold detection, and resilient data persistence.

\subsection{Logging and Persistence}
Metrics are persisted locally in CSV format for reliability and traceability.

\subsection{Webhook Integration}
\subsubsection{n8n Webhook Client}
HTTP POST requests send JSON metrics every 5 seconds with retry logic.

\subsubsection{Webhook Payload}
The payload includes 17 metrics: timestamp, queue count, rates, waiting time, and uncertainty level.

\subsection{Threshold Detection}
A dedicated module flags critical situations:
\begin{itemize}
    \item excessive waiting time,
    \item high queue backlog,
    \item arrival spikes,
    \item service degradation,
    \item instability or high uncertainty.
\end{itemize}

\section{Code Organization}
The project is organized into specialized modules:
\begin{itemize}
    \item \texttt{detector.py}: people detection,
    \item \texttt{tracker.py}: object tracking,
    \item \texttt{zone\_manager.py}: polygon-zone management,
    \item \texttt{queue\_analyzer.py}: queue metric computation,
    \item \texttt{uncertainty.py}: Bayesian uncertainty estimation,
    \item \texttt{threshold\_detector.py}: critical condition detection,
    \item \texttt{gui/app.py}: graphical interface and multi-video launcher.
\end{itemize}
\label{sec:code_org}

\section{Software Validation}
Unit tests are included (notably for the interface and uncertainty modules).
Validation also covered real-time runtime scenarios on video streams.

\section{Challenges and Solutions}
This project required several engineering adjustments:
\begin{itemize}
    \item stabilization of GUI rendering,
    \item improved robustness of zone selection,
    \item calibration of the performance/accuracy trade-off,
    \item handling edge cases in the queueing model.
\end{itemize}
